The 2018 paper \citet{corcoran2018spatial} introduced a Hilbert-curve-based mapping that converted voxelized 3D structures into 2D and 1D arrays so that lower-dimensional convolutional networks could learn from them efficiently. That framing was timely in the volumetric-CNN era, but it no longer captures the strongest scientific value of the idea. In this paper, we revisit the original work through the lens of modern 3D and biomolecular learning and argue for a different conclusion: serialization is not merely a workaround for expensive volumetric processing, but a first-class interface between irregular 3D structure and sequence-native model families. We support this reinterpretation by synthesizing evidence from point attention, serialized point transformers, state-space models, neural surface reconstruction, and modern biomolecular representation learning \citep{zhao2021pointtransformer, wu2024ptv3, liang2024pointmamba, zhang2024voxelmamba, li2025noksr, park2025atomae}. The result is a distinct perspective manuscript with three concrete contributions: a cross-era taxonomy of serialized structural learning, a design framework organized around locality, token efficiency, channel semantics, and reversibility, and a publication agenda for biomolecular AI. We do not claim new benchmark results. Instead, the contribution is conceptual: the original Hilbert-mapping paper is repositioned as an early precursor to a now-active design space in which serialization shapes what modern models can represent, scale, and explain.
